\documentclass[11pt]{article}
\usepackage[a4paper,margin=27mm]{geometry}
\usepackage{amsmath,amssymb,amsthm}
\newtheorem{lemma}{Lemma}
\newtheorem{proposition}[lemma]{Proposition}
\title{Reusing the original unit consequences in a proof check}
\author{Human Compression project}
\date{}
\begin{document}\maketitle
Let $F$ be the original propositional formula in conjunctive normal form.
The historical RUP checker keeps every original clause and every previously
verified addition. It ignores deletion instructions after parsing them.
For an addition $C$, it requires unit propagation to derive a contradiction
from the retained clauses and the negations of the literals of $C$.
These are the conventions of the existing checker in this project.

Repeatedly deriving the same original unit consequences is unnecessary.
The following construction implements the fixed-assignment reduction
described in the retained note \emph{Fixed original-unit reduction for the
historical RUP checker}. The propagation algorithm and clause parser are
imported unchanged from the accepted checker. This is a computational
improvement of that verification method; no new graph theorem or claim of
priority for unit propagation is made.

Run unit propagation on $F$, with no assumptions. If it does not find a
contradiction, let $R$ be its complete assignment. Every literal in $R$ has
a unit derivation from $F$. The assignment $R$ is frozen for the remainder
of the check. In particular, values derived while checking a proposed
addition are never inserted into $R$.

For a clause $C$, let $C\vert_R$ be true if $C$ has a true literal under
$R$. Otherwise obtain $C\vert_R$ by removing its false literals. Apply
this operation to a formula by dropping its true clauses.

\begin{lemma}
Suppose that $H$ contains $F$ and that $R$ is the consistent unit closure
of $F$. A clause $C$ passes the historical RUP test on $H$ if and only if
either $C\vert_R$ is true, or $C\vert_R$ passes that test on $H\vert_R$.
\end{lemma}
\begin{proof}
If $C\vert_R$ is true, negating $C$ contradicts a literal with a unit
derivation from $F$. This gives a unit contradiction from $H$ and those
negations.

Otherwise the negations of the removed literals already hold in $R$.
Start propagation from $H$ and the negation of $C$. Each original derivation
of a literal in $R$ can still be followed: either the same literal is
derived, or an earlier contradiction occurs. If no such contradiction
occurs, all of $R$ has been obtained. After this stage a clause already
satisfied by $R$ has no further effect, and all other clauses and remaining
assumptions are exactly their residuals. Every subsequent unit step therefore
is a unit step in the residual problem. Conversely, every residual step
lifts to an original step after the derivations of $R$.

Exhaustive unit propagation cannot depend on the order of its steps as to
whether a contradiction is found. Indeed, a finite unit derivation can be
followed in any exhaustive propagation: each required literal is already
present, is eventually forced, or conflicts with an earlier assignment.
Apply this observation to the two corresponding derivations just described.
\end{proof}

\begin{proposition}
The residual checker admits exactly the same additions as the historical
checker, including after arbitrary deletion records. Acceptance of an
explicit empty addition proves that $F$ is unsatisfiable.
\end{proposition}
\begin{proof}
Maintain in the residual checker the residuals of $F$ and of all previously
admitted additions. A root-satisfied addition is already entailed by $F$
and has no residual clause to store. For every other addition, the lemma
identifies the two admission tests. Add its residual only after the test
passes. Induction on the addition sequence proves the assertion. Both
historical checks ignore deletions, so these do not change the induction.

Every RUP addition is a consequence of the retained original clauses and
previous additions: a truth assignment satisfying them and falsifying the
addition would satisfy the negated assumptions and could not yield a unit
contradiction. A second induction makes every admitted addition a consequence
of $F$. The explicit empty clause is consequently impossible in a model
of $F$.

If the initial propagation already finds a contradiction, $F$ itself is
unit-inconsistent. Every subsequent addition passes its RUP test, and the
same conclusion follows. The implementation still parses the entire proof
and requires an explicit empty addition.
\end{proof}

The implementation retains the original variable numbering and input-clause
count. It parses every input and proof clause before taking a shortcut,
including deletion instructions and lines following an empty addition.
A nonempty original clause whose residual is empty does not satisfy the
explicit-empty requirement. The assignment used for residualization comes
only from a separate propagator containing $F$; the residual propagator has
its own mutable assignments and watched clauses.

The test harness reuses the existing independent set-propagation oracle and
compares the first invalid addition as well as the final result. It includes
the retained 317 cases, ten additional cases for fixed assignments and
temporary assumptions, and six malformed proof records. The added cases
include different original clauses with the same residual, a root-false
candidate, and a nonempty clause that cannot serve as the concluding empty
addition. These checks support the implementation; the preceding induction
is the mathematical justification. Application to any saved graph proof
still requires its exact formula reconstruction and an actual successful
proof check.
\end{document}
